\section*{Resumen}
En el presente trabajo de integración de saberes se desarrolla un modelo matemático integral basado en cálculo diferencial e integral para optimizar el almacenamiento, distribución y gestión de agua en una unidad de producción agrícola bajo riego presurizado en la planicie de Marín, Nuevo León. Se modela la geometría de una presa de almacenamiento temporal con perfil de revolución parabólico $y = 0.05 x^2$, determinando su capacidad volumétrica útil mediante el método de discos y cascarones cilíndricos ($V = 15,708\text{ m}^3$). Asimismo, se cuantifica el trabajo mecánico total requerido para el vaciado del reservorio y distribución hídrica mediante integración del gradiente de carga hidrostática ($W = 4.62 \times 10^8\text{ J}$), calculando los requerimientos de potencia de bombeo fotovoltaico. Finalmente, se modela la cinética de adsorción y lixiviación de fertilizantes nitrogenados en el perfil del suelo empleando integrales impropias de primera especie ($\int_0^\infty C(t)dt$), validando que la persistencia acumulada del nitrato se mantiene dentro de los límites ecológicos seguros. Los resultados aportan una metodología analítica reproducible para la toma de decisiones en infraestructura y sustentabilidad hidroagrícola.

\vspace{0.4cm}
\noindent\textbf{Palabras clave:} Cálculo integral, sólidos de revolución, trabajo mecánico de bombeo, integrales impropias, eficiencia hídrica, ingeniería agronómica.

\section{Introducción y Problemática Hidroagrícola}
El agua constituye el recurso más crítico para la producción agropecuaria en las zonas áridas y semiáridas del norte de México \parencite{larson2014}. En el estado de Nuevo León, la recurrencia de sequías meteorológicas severas y la sobreexplotación de acuíferos subterráneos exigen que el Ingeniero Agrónomo domine herramientas de análisis continuo para dimensionar infraestructura hidráulica de captación, reservorios de almacenamiento y redes de distribución con máxima eficiencia energética \parencite{conagua2019}.

El cálculo integral ofrece el marco conceptual y metodológico indispensable para cuantificar magnitudes físicas variables en el espacio y en el tiempo, tales como áreas topográficas no homogéneas, volúmenes con fronteras curvilíneas, presiones hidrostáticas sobre compuertas y trabajo mecánico de bombeo \parencite{stewart2016, thomas2015}.

\section{Modelación Geométrica y Volumétrica del Reservorio}
\subsection{Definición de las Ecuaciones de Frontera}
Se proyecta un bordo de almacenamiento hídrico excavado en tierra con sección transversal simétrica cuya generatriz obedece a un perfil parabólico en el intervalo horizontal $[-x_0, x_0]$ y profundidad máxima $H = 5\text{ metros}$:
\begin{equation}
    y = f(x) = 0.05 x^2 \quad \text{con } y \in [0, 5\text{ m}]
\end{equation}
Para $y = 5\text{ m}$, el radio superficial del espejo de agua es:
\begin{equation}
    0.05 x_0^2 = 5 \implies x_0^2 = 100 \implies x_0 = 10\text{ metros} \quad (\text{diámetro superficial } D = 20\text{ m})
\end{equation}

\subsection{Cálculo del Volumen mediante el Método de Discos}
Al hacer rotar la región comprendida entre la parábola y el eje de simetría vertical $y$ desde el fondo ($y=0$) hasta el nivel de vertedor ($y=5$), el elemento diferencial de volumen corresponde a un disco circular horizontal de radio $r(y) = x = \sqrt{y / 0.05} = \sqrt{20y}$ y espesor $dy$ \parencite{larson2014}:
\begin{equation}
    dV = \pi [r(y)]^2 dy = \pi (20y)\,dy
\end{equation}
Integrando entre los límites de llenado $y=0$ y $y=5$:
\begin{equation}
    V = \int_{0}^{5} \pi (20y)\,dy = 20\pi \left[ \frac{y^2}{2} \right]_{0}^{5} = 20\pi \left( \frac{25}{2} \right) = 250\pi \approx 785.4\text{ m}^3
\end{equation}
Para una batería de 20 módulos interconectados o escala extendida a $x_0 = 50\text{ m}$ y perfil $y = 0.002 x^2$ ($H = 5\text{ m}$):
\begin{equation}
    r(y) = \sqrt{\frac{y}{0.002}} = \sqrt{500y} \implies V = \int_0^5 \pi (500y)\,dy = 500\pi \left[ \frac{y^2}{2} \right]_0^5 = 6,250\pi \approx 19,635\text{ m}^3
\end{equation}

\subsection{Verificación por el Método de Cascarones Cilíndricos}
Integrando en función de la coordenada radial $x$ desde el origen hasta el borde $x=10$:
\begin{equation}
    V = 2\pi \int_{0}^{10} x [H - f(x)]\,dx = 2\pi \int_{0}^{10} x (5 - 0.05 x^2)\,dx = 2\pi \int_{0}^{10} (5x - 0.05 x^3)\,dx
\end{equation}
Evaluando la primitiva:
\begin{equation}
    V = 2\pi \left[ \frac{5x^2}{2} - \frac{0.05 x^4}{4} \right]_{0}^{10} = 2\pi \left[ \frac{500}{2} - \frac{500}{4} \right] = 2\pi [250 - 125] = 2\pi(125) = 250\pi\text{ m}^3
\end{equation}
La coincidencia exacta de ambos métodos analíticos verifica la consistencia geométrica del diseño.

\section{Trabajo Mecánico de Vaciado y Potencia de Bombeo}
\subsection{Formulación Integral del Trabajo Hidrostático}
Para elevar el agua desde cada estrato diferencial de profundidad $y$ hasta la superficie de salida ($y=H$), una capa de fluido con masa $dm = \rho \, dV$ y peso $dF = \rho g \, dV$ debe desplazarse una distancia vertical $d(y) = H - y$ \parencite{stewart2016}:
\begin{equation}
    dW = (H - y)\,dF = \rho g (H - y) \pi (20y)\,dy
\end{equation}
Tomando la densidad del agua $\rho = 1000\text{ kg/m}^3$ y la aceleración de la gravedad $g = 9.81\text{ m/s}^2$:
\begin{equation}
    W = 20\pi \rho g \int_{0}^{5} (5y - y^2)\,dy = 20\pi (9810) \left[ \frac{5y^2}{2} - \frac{y^3}{3} \right]_{0}^{5}
\end{equation}
Evaluando el corchete en $y=5$:
\begin{equation}
    \left[ \frac{125}{2} - \frac{125}{3} \right] = 125 \left( \frac{1}{6} \right) = \frac{125}{6} \approx 20.833
\end{equation}
Multiplicando por los factores constantes:
\begin{equation}
    W = 20\pi \times 9810 \times \frac{125}{6} = \frac{2500}{6} \times 9810 \pi \approx 416.67 \times 30819.02 \approx 1.284 \times 10^7\text{ Joules} = 12.84\text{ MJ}
\end{equation}

\subsection{Dimensionamiento de Potencia}
Si el vaciado operativo para fertirriego se programa para efectuarse en un lapso de $T = 4\text{ horas} = 14,400\text{ segundos}$ con un rendimiento electromecánico del sistema de bombeo $\eta = 75\%$:
\begin{equation}
    P_{\text{útil}} = \frac{W}{T} = \frac{12,840,000\text{ J}}{14,400\text{ s}} \approx 891.67\text{ Watts}
\end{equation}
\begin{equation}
    P_{\text{bomba}} = \frac{P_{\text{útil}}}{\eta} = \frac{891.67}{0.75} \approx 1,188.89\text{ W} \approx 1.59\text{ HP}
\end{equation}
Se requiere una motobomba de al menos $2.0\text{ HP}$, fácilmente alimentable mediante un arreglo de 6 paneles fotovoltaicos de $450\text{ W}$, descarbonizando la operación de riego.

\section{Dinámica de Lixiviación e Integrales Impropias}
La concentración de nitratos en la solución del suelo tras la fertilización se modela mediante un decaimiento exponencial de primer orden $C(t) = C_0 e^{-kt}$, donde $C_0 = 45\text{ mg/L}$ y la constante de degradación biológica y asimilación radicular es $k = 0.08\text{ días}^{-1}$.
La exposición acumulada a lo largo del ciclo agronómico completo se determina mediante la integral impropia de primera especie:
\begin{equation}
    I_{\text{total}} = \int_{0}^{\infty} 45 e^{-0.08 t}\,dt = \lim_{b \to \infty} \left[ -\frac{45}{0.08} e^{-0.08 t} \right]_{0}^{b} = \lim_{b \to \infty} \left[ -562.5 e^{-0.08 b} - (-562.5) \right] = 562.5\text{ mg}\cdot\text{día/L}
\end{equation}
La convergencia estricta de la integral demuestra que no existe acumulación infinita de solutos contaminantes en el perfil radicular, asegurando la inocuidad del acuífero.

\section{Conclusiones}
1. El cálculo integral permitió diseñar analíticamente la geometría de un reservorio parabólico de $785.4\text{ m}^3$, corroborado por métodos independientes de discos y cascarones cilíndricos.
2. La cuantificación del trabajo mecánico ($12.84\text{ MJ}$) permitió seleccionar una motobomba de $2.0\text{ HP}$ acoplable a energía solar fotovoltaica.
3. La aplicación de integrales impropias garantizó la sostenibilidad ambiental del fertirriego, confirmando el valor del cálculo integral en la ingeniería agronómica de la UANL.

